\subsection{A complete countable presentation for components in a stable range}

The following computation specifies an abelian group by generators and
all its relations. It does not assert finite generation, a finite
invariant-factor decomposition, or an algorithm deciding equality of
arbitrary words. In particular it supplies the extension law that is
not determined by the outer terms of the Hatcher--Wagoner--Igusa exact
sequence.

\begin{saresult}{A.119}{Proved under stated hypotheses}{An explicit countable component presentation}
Let $B$ be a closed definable manifold whose finite PL model $K$ is
smoothable and has dimension at least $9$ on every component. The
construction below, applied to each connected component of $K$, gives
a countable abelian presentation $E(K)$ and an isomorphism
\[
                   \pi_0P_{\rm def}(B)\cong E(K).
\]
The group operation on the right is addition of integral cycles. The
presentation depends only on the finite simplicial homotopy model $K$,
and is valid over every real closed field and every fixed o-minimal
expansion. It also gives the stable smooth concordance component group
for any compact smooth manifold, without a dimension restriction.
It makes no unstable assertion for a nonsmoothable PL manifold or for
an arbitrary manifold of dimension at most $8$.
\end{saresult}

\paragraph{The presentation.}
We first give the construction for connected $K$. Collapse a maximal
tree in its $1$-skeleton and use the resulting reduced finite simplicial
set. Let $G$ be its Kan loop group and put $R=\Sigma^\infty_+G$, the
spherical group ring. A countable simplicial model for these objects
and all the operations used below is specified in the proof.
Let
\[
 H_n=\operatorname{GL}_n(R),\qquad
 W_n=G^n\rtimes\Sigma_n.
\]
Here $H_n$ is the derived simplicial monoid of homotopy automorphisms
of the free right $R$-module $R^{\vee n}$; it is not the discrete
group $\operatorname{GL}_n(\Z[\pi_1K])$. Permutations of summands and
left multiplication by elements of $G$ define $W_n\to H_n$.
Stabilize by adjoining an identity summand and form
\[
       C=\operatorname*{hocolim}_n BH_n,\qquad
       D=\operatorname*{hocolim}_n BW_n,
       \qquad f:D\longrightarrow C.                 \tag{A.119.1}
\]
We use countable based simplicial telescope models. If necessary,
collapse a maximal tree to give each model a single vertex; the map
$f$ can equivalently be retained before this collapse, using the
based lift described below. All chains in what follows are normalized
simplicial chains with integral coefficients and finite support.

For a connected one-vertex simplicial set $Y$, put
\[
 L_Y=C_1(Y)/\partial C_2(Y)=H_1(Y;\Z).
\]
Its universal abelian covering $Y^{\rm ab}$ has $q$-simplices
$(a,\sigma)$ with $a\in L_Y$ and $\sigma\in Y_q$. The first face is
\[
 d_0(a,\sigma)=
 (a+[\sigma|[0,1]],d_0\sigma),
\]
and the other faces and all degeneracies retain $a$. For $q=0$ there
are no face maps. The triangle relations in $L_Y$ verify the
simplicial identities. For a model with several vertices the same
formula uses the edge labels obtained by joining the two endpoints
to the base vertex in the chosen maximal tree. This specifies a
connected cover corresponding to
$\ker(\pi_1Y\to H_1Y)$. The based map $f$ induces a based map
$f^{\rm ab}:D^{\rm ab}\to C^{\rm ab}$.

Write $Z_2(Y)=\ker(\partial:C_2(Y)\to C_1(Y))$. The generators of
$E(K)$ are symbols $[z]$ for every $z\in Z_2(C^{\rm ab})$. Its
complete list of relations is
\[
\begin{aligned}
 [z+z']&=[z]+[z'] &&(z,z'\in Z_2(C^{\rm ab})),\\
 [\partial c]&=0 &&(c\in C_3(C^{\rm ab})),\\
 [(f^{\rm ab})_*w]&=0 &&(w\in Z_2(D^{\rm ab})).
\end{aligned}                                                   \tag{A.119.2}
\]
Thus it is also the explicit quotient
\[
 E(K)=\frac{Z_2(C^{\rm ab})}
 {\partial C_3(C^{\rm ab})+(f^{\rm ab})_*Z_2(D^{\rm ab})}.
                                                               \tag{A.119.3}
\]
All faces in these formulas are ordinary bar, simplicial, and
telescope faces; in particular every boundary is the alternating sum
$\partial\sigma=\sum_i(-1)^id_i\sigma$. For a disconnected $K$
take the finite product of the presentations for its components.

\begin{proof}
\emph{Countable models determined by finite data.}
For the reduced simplicial set $X$ obtained from $K$, the Kan loop
group is explicitly
\[
 G(X)_q=F(X_{q+1})/\langle[s_0x]=1:x\in X_q\rangle,
\]
with the following homomorphisms on generators:
\[
 d_0[x]=[d_1x][d_0x]^{-1},\qquad
 d_i[x]=[d_{i+1}x]\ (i>0),\qquad
 s_i[x]=[s_{i+1}x].
\]
These are the convention of \cite[V.5, (5.1), p.~276]
{GoerssJardine1999}; their Theorem~5.10 and Corollary~5.11
(p.~285) give the loop-space comparison. All elements are finite
words in the given simplices, so the simplicial sets are countable.
Use simplicial symmetric spectra for $R$ and their usual
stable module model structure. Here are precise choices that avoid
using an uncountable singular mapping space. Let $\mathcal F_R$
be the category of finite cell right $R$-modules, with their cell
presentations retained. Cells are the free modules on
$F_j\partial\Delta^q_+\to F_j\Delta^q_+$. All attaching maps and all
maps between such objects are specified by the images of finitely
many cell generators in countable simplicial sets. Consequently the
category has a countable skeleton and countable morphism sets.

For acyclic attachments use the countable set $R\wedge J$, where
$J=FI_\Lambda\cup\{c_j\mathbin\square i\}$ is the explicit set of
\cite[Definition~3.4.9, pp.~181--182; Theorem~3.4.4, p.~180]
{HoveyShipleySmith2000}. Here $i$ runs over boundary inclusions,
$c_j:F_{j+1}S^1\to C\lambda_j$ is the mapping-cylinder
cofibration of $\lambda_j:F_{j+1}S^1\to F_jS^0$, and
$FI_\Lambda$ consists of the free horn inclusions. Their domains
and codomains are finite free-cell spectra. Smashing with $R$
gives finite-cell module maps. The transferred module model
structure is \cite[Theorem~4.1(1), pp.~498--499; proof p.~504]
{SchwedeShipley2000}, applied to $R^{\rm op}$ for right modules.
We also retain the finite cylinder inclusions
$P\to P\wedge\Delta^1_+$.

Take finite pushouts and composites of these acyclic attachments,
and formally invert them. This supplies a finite-certificate description
of every stable equivalence between finite cell modules. To check the
last assertion, construct a fibrant replacement by the small-object
argument using the stated generators, in countably many stages. Any
map from a finite cell module, or from its tensor with a finite
simplicial parameter set, into that replacement factors through a
finite relative cell extension: it uses finitely many generators,
and each attaching map uses finitely many earlier generators.
Induction through the finitely many stages involved gives a finite
set closed under those dependencies. If a map between finite cell
modules is a stable equivalence, its inverse in the model homotopy
category and the two inverse homotopies have precisely such finite
factorizations. They therefore give an inverse after formally
inverting finite acyclic extensions. Conversely every map inverted
by this procedure is a stable equivalence, because all the inverted
attachments are such and cylinder homotopies preserve them. Thus
no geometric concordance test, and no unspecified extra invariant,
is used to select these maps.

Use the hammock localization of this countable category and its
stable equivalences. Its simplices are finite diagrams, so it remains
countable. It computes the derived module mapping spaces:
\cite[Proposition~4.4 and Application~4.6, p.~432]
{DwyerKan1980Functions} gives the function-complex comparison, and
the preceding finite-factorization argument restricts that comparison
to the finite cell objects. It applies also to any finite simplicial
parameter diagram, so verifies both surjectivity and injectivity on
all homotopy groups of the mapping spaces. Select the invertible
components of the endomorphism monoid of $R^{\vee n}$ to obtain
$H_n$. Invertibility itself is witnessed by a finite inverse diagram
and two finite paths to the identity. Direct sum preserves all the
specified attachments and gives the direct-sum coherences on these
models. They can be kept as coherent operations, or rectified by the
simplicial bar construction; that construction uses finite strings
and preserves countability. This specifies the models in (A.119.1)
from finite $K$. It does not assert a terminating procedure for
recognizing a non-equivalence or for solving the resulting group
word problem.

\emph{Acyclic group completion and the two covers.}
Under direct sum the spaces
\[
 \mathcal M=\coprod_{n\geq0}BH_n,
 \qquad \mathcal N=\coprod_{n\geq0}BW_n
\]
are homotopy commutative monoids. The homotopy commutativity is
given by block-permutation conjugation. It is this assertion about
$\mathcal M,\mathcal N$, not about the uncompleted connected space
$C$, that is used here. Their component monoids are $\mathbb N$, and
successively adding one identity summand is cofinal. The
permutation-matrix models in \cite[\S2, p.~247]{Jahren2010}
identify their group completions as
\[
 \Omega B\mathcal M\simeq A(|K|),
 \qquad \Omega B\mathcal N\simeq Q(|K|_+),
                                                               \tag{A.119.4}
\]
and identify their induced map with the unit
$\iota:Q(|K|_+)\to A(|K|)$. Only the indicated rank-zero
components matter; no identification of all projective $K_0$
classes with free rank is being made.

By \cite[Proposition~1.1 and proof pp.~1--7, especially p.~3]
{Gritschacher2026Acyclicity}, the telescope-to-group-completion maps
are acyclic. Its hypotheses hold since the component monoid is
cofinal and central in the homology ring of a homotopy commutative
monoid. Thus the maps from $C$ and $D$ to the corresponding
rank-zero components in (A.119.4) induce abelianization on fundamental
groups and homology isomorphisms with every local coefficient
system. Pull back the simply connected covers of those components.
The pullbacks are exactly $C^{\rm ab}$ and $D^{\rm ab}$, and the
lifted maps are integral homology isomorphisms. Hurewicz on the
simply connected targets gives natural isomorphisms
\[
 H_2(C^{\rm ab};\Z)\cong\pi_2A(|K|),
 \qquad H_2(D^{\rm ab};\Z)\cong\pi_2Q(|K|_+).
                                                               \tag{A.119.5}
\]
The square involving $f^{\rm ab}$ and the lifted unit commutes up
to the specified based homotopy, so these isomorphisms respect the
homomorphism whose cokernel appears in (A.119.3).

\emph{The whole component group, including the extension law.}
The split unit and the stable parametrized h-cobordism theorem give
\[
 A(|K|)\simeq Q(|K|_+)\times\operatorname{Wh}^{\rm Diff}(|K|),
 \qquad
 \pi_0P_{\rm Diff}^{\rm st}(B)
       \cong\pi_2\operatorname{Wh}^{\rm Diff}(|K|).
\]
The split sequence and its unit are read in
\cite[p.~247 and the basic diagram p.~248]{Jahren2010}; the full
geometric comparison is \cite[Theorem~0.1, p.~2; Definition~1.3.4
and (1.3.5), p.~18]{WaldhausenJahrenRognes2013}.
Here the compatibility with the original composition operation
requires a delooping argument. For any simplicial group $P$, the
map $\pi_0P\to\pi_1BP$ takes $g$ to its bar edge; the bar
$2$-simplex $(g,h)$ identifies the edge of the product with the
concatenated edges, using the matching bar convention. Thus this
map is an isomorphism of groups. By
\cite[Remark~1.1.11 and (1.1.12), p.~12]
{WaldhausenJahrenRognes2013}, the component $H(M)_0$ of the
trivial h-cobordism is $BP_{\rm Diff}(M)$. More precisely, the
side-fixed concordance group includes by a group homomorphism and
weak equivalence into the automorphism group of $M\times I$
relative to $M\times0$, whose bar construction classifies that
component. Stabilization on h-cobordisms is the functor
$W\mapsto W\times J$, taking each automorphism $g$ to
$g\times\mathrm{id}_J$ and preserving composition. Hence these
group identifications commute with stabilization. The stabilization
maps of $H$ are simplicial cofibrations
\cite[Definition~1.1.3, p.~8]{WaldhausenJahrenRognes2013}, so
their fundamental groups commute with the colimit. Applying
$\pi_1$ to the based stable h-cobordism equivalence gives
\[
 \operatorname*{colim}_j\pi_0P_{\rm Diff}(M\times J^j)
 \cong\pi_1H^{\rm st}(M)_0
 \cong\pi_1\Omega\operatorname{Wh}^{\rm Diff}(M)
 =\pi_2\operatorname{Wh}^{\rm Diff}(M)
\]
as groups. A based equivalence induces a homomorphism on
$\pi_1$; no preservation of an unlooped multiplication is being
assumed. The split-unit maps and Hurewicz maps are also
homomorphisms. It follows from (A.119.5) that the latter group,
with this precise operation, is exactly (A.119.3), where Hurewicz
carries it to addition of cycles.
The definition (A.119.2) therefore includes all relations between
lifts of quotient classes. It leaves no separate splitting or
extension class to be chosen. Invariance under a change of the
simplicial or spectral models follows from these natural
identifications, rather than from a claim that their sets of
simplices are identical.

\emph{The stated unstable and definable range.}
For a smoothable $n$-dimensional PL manifold choose a smoothing.
Burghelea--Lashof's component isomorphism
\cite[Theorem~6.1, p.~41; proof pp.~42--43]
{BurgheleaLashof1974II}, already checked in A.103, identifies smooth
and PL concordance components. By
\cite[Theorem~1.4.1, p.~21]{WaldhausenJahrenRognes2013}, smooth
stabilization is $k$-connected when
$n\geq\max\{2k+7,3k+4\}$. With $k=1$, every stabilization from
$n\geq9$ onward is $1$-connected, hence bijective on components.
Consequently the unstable group in this range equals the stable
group just presented. A.113 identifies its components with the
definable group over every real closed field. The finite product
statement follows because a concordance fixed on the bottom cannot
permute connected components.

Smoothability is essential to this particular range deduction:
the opened book explicitly separates the general PL stability
problem from its smoothable Corollary~1.4.2 (p.~22). No unstable
conclusion outside the stated class is obtained by this argument.
\end{proof}
